\subsection{理论中的代入}

设 \(\mathfrak{c}\) 为一个理论，令 \(A_{1}, A_{2}, \ldots, A_{n}\) 为其显式公理，令 \(\pmb{x}\) 为一个字母，并且令 \(\pmb{T}\) 为 \(\mathfrak{c}\) 的一项。设 \((\pmb{T}|\pmb{x})\mathfrak{c}\) 为这样的理论，其符号和方案与 \(\mathfrak{c}\) 的相同，且其显式公理为 \((\pmb{T}|\pmb{x})\pmb{A}_{1}, (\pmb{T}|\pmb{x})\pmb{A}_{2}, \ldots, (\pmb{T}|\pmb{x})\pmb{A}_{n}\) .

C2。设A是理论 \(\mathfrak{G}\) 中的一个定理，设T是 \(\mathfrak{G}\) 的一个项，并设x是一个字母。则 \((T|x)A\) 是理论 \((T|x)\mathfrak{G}\) 中的定理。

设 \(R_1, R_2, \ldots, R_n\) 是 \(\mathfrak{G}\) 中的一个证明，在其中 \(A\) 出现。考虑序列 \((T|x)R_1, (T|x)R_2, \ldots, (T|x)R_n\) ，这是 \(\mathfrak{G}\) 中的关系序列，依据是CF8（§1，第4号）。我们将证明这个序列是该理论 \((T|x)\mathfrak{G}\) 中的一个证明；这将确立该准则。如果 \(R_k\) 是 \(\mathfrak{G}\) 的一个隐式公理，那么 \((T|x)R_k\) 同样是 \(\mathfrak{G}\) 的一个隐式公理（编号1），因此也是 \((T/x)\mathfrak{G}\) 的。如果 \(R_k\) 是 \(\mathfrak{G}\) 的一条显式公理，那么 \((T|x)R_k\) 是 \((T|x)\mathfrak{G}\) 的一条显式公理。最后，如果 \(R_k\) 前面有关系 \(R_i\) 和 \(R_j\) ，其中 \(R_j\) 是 \(R_i \Longrightarrow R_k\) ，那么 \((T|x)R_k\) 前面是 \((T|x)R_i\) 和 \((T|x)R_j\) ，后者与 \((T|x)R_i \Longrightarrow (T|x)R_k\) （准则CS5）相同。

C3. 设 A 是理论 G 的一个定理，设 T 是 G 的一个项，并设 x 是一个不是 G 的常量的字母。则 \((T|x)A\) 是 G 的一个定理。

这直接从 C2 推出，因为 x 不出现在 \(\mathfrak{G}\)的显式公理中。更特别地，如果 \(\mathfrak{G}\) 不包含任何显式公理，或者显式公理中不包含任何字母，那么准则C3对字母x的应用不受限制。

